\documentclass[11pt,a4paper]{article}
\usepackage[T1]{fontenc}
\usepackage[utf8]{inputenc}
\usepackage[margin=22mm]{geometry}
\usepackage{lmodern}
\usepackage{amsmath,amssymb}
\usepackage{enumitem}
\usepackage{xcolor}
\usepackage{microtype}
\usepackage{hyperref}
\definecolor{epflred}{HTML}{E3000B}
\hypersetup{colorlinks=true,urlcolor=epflred,linkcolor=epflred,
  pdftitle={ENG-654 Project 5: From Positioning-Arm Geometry to 6R Cuspidality}}
\setlength{\parindent}{0pt}
\setlength{\parskip}{6pt}
\setlist[itemize]{leftmargin=*,itemsep=5pt,topsep=4pt}
\urlstyle{tt}

\begin{document}
{\small\textbf{EPFL \textbar\ ENG-654}\hfill Kinematics-Grounded Motion Planning for Robots}
\vspace{5mm}

{\color{epflred}\Large\bfseries Project 5}

{\LARGE\bfseries From Positioning-Arm Geometry to 6R Cuspidality\par}

\textbf{Format:} Group of two students; simulation-based project.\\
\textbf{Prerequisites:} D--H modelling, spherical-wrist inverse kinematics,
Jacobians, singularities, and configuration-space connectivity.

\section*{Motivation}
Attaching a spherical wrist turns a position manipulator into a full-pose
robot, but it does not erase the positioning arm's global structure. A useful
classification must explain both the arm's possible posture changes and the
wrist's orientation singularities. This project studies the relationship rather
than assuming that every wrist-partitioned robot behaves like PUMA.

\section*{Task}
\textbf{Overall objectives.} Analyze a selected family of wrist-partitioned 6R robots formed from
noncuspidal and cuspidal 3R arms, and relate full-pose IK multiplicity and
connectivity to the two kinematic modules.

\begin{itemize}
  \item \textbf{Specify and validate the family.} Use the two idealized 3R arms
$a=(a_1,2,1.5)$, $d=(1,1,0)$, $\alpha=(\pi/2,\pi/2,0)$, with $a_1=0$
and $a_1=1$ for the noncuspidal and cuspidal members respectively. State a
common length unit. Attach the same ideal spherical wrist with an explicit fixed
tool transform. Define the wrist rotation convention and verify its axis
concurrency. Use the supplied custom 6R robot as a concrete example of a
3R positioning arm plus spherical wrist; retain the stated idealized arm
dimensions for the parameter study.

  \item \textbf{Factor the Jacobian geometrically.} Express the geometric Jacobian at
the wrist center in a common frame and show the block structure giving
$\det J=\det J_{A,v}\det J_{W,\omega}$. Interpret arm and wrist rank loss;
confirm that changing the point at which the twist is expressed preserves rank.

  \item \textbf{Map the positioning-arm structure.} Plot $\det J_{A,v}$ and its zeros
in $(q_2,q_3)$, overlay arm IKs, and map the two-/four-arm-IK regions in
$(\rho,z)$. For fixed regular tool orientations, count the resulting full-pose
solutions and explain when each arm solution admits two wrist flips.

  \item \textbf{Test the inheritance of connectivity.} Use one supplied candidate
nonsingular arm connection for the cuspidal member and solve the wrist along
it while maintaining a selected orientation. Monitor both Jacobian factors.
Explain why a failed fixed-orientation lift does not by itself refute
cuspidality of the unrestricted full-pose robot.

  \item \textbf{Compare structural conclusions.} Relate singular separators and IK
counts for both members. Distinguish global classification of the ideal
unconstrained robot from the feasibility of a particular orientation-constrained
path. Repeat one target test after changing the fixed tool offset, explaining
the required pose transformation.
\end{itemize}

\newpage
\section*{Specifics and scope}
Limit the study to two arm geometries, one common spherical-wrist
architecture, three regular targets per member, and one candidate connecting
path. The instructor will provide the arm IK routine and candidate connection.
Derive the wrist solution and Jacobian factorization; a classification of all
6R robots is not required. Apply joint limits only as a separately reported
feasibility filter.

Check both translation and orientation residuals
for every claimed full-pose IK. Document angle periodicity and limits,
Jacobian conventions, and any scaling used for singular values. State all
fixed coordinates in reduced plots; a two-dimensional slice does not show
the entire 6R singular set. Refine decisive transitions, and distinguish
sampled evidence from a continuous-path guarantee. Collision freedom is
not implied by these kinematic checks.

\section*{Expected outcome}
Explain how arm geometry, wrist flips, and their singularities organize
the full-pose inverse problem. Deliver paired arm maps, full-pose IK counts,
the factorization, and a documented comparison of the selected path's arm and
full-robot regularity.

Submit reproducible code, model parameters and test data, the required plots,
a concise 5--6-page report, and a short simulation demonstration. Explain
both successful cases and failures; conclusions must follow the numerical
and geometric evidence rather than an expected outcome.

\section*{Provided materials and implementation}
The custom 6R robot description (.urdf) and link meshes (.stl) will be provided
to the students. Its positioning arm is orthogonal with non-intersecting
consecutive axes, followed by a spherical wrist. Inspect wrist-axis concurrency
and base/tool transforms using a robotics library or your own visualization.
The arm IK routine and candidate nonsingular arm connection will also be
provided for the stated idealized geometries.

Use standard D--H transforms
$T_{i-1,i}=R_z(q_i)T_z(d_i)T_x(a_i)R_x(\alpha_i)$, with zero joint-angle
offsets for the idealized arms. Their dimensions must not be inferred from
the supplied meshes. Derive the wrist orientation solution and full Jacobian
in the same frame convention, then verify all full-pose IKs by forward
kinematics. Use plotting tools to display singularities and animate the
selected connection.

\end{document}
